\chapter{Conclusioni e Sviluppi Futuri}

Questa tesi ha affrontato una questione ancora poco esplorata nel campo del machine unlearning: il ruolo della variabilità linguistica nell'efficacia dei processi di rimozione selettiva della conoscenza. Mentre la maggior parte degli approcci esistenti si concentra sulla progettazione di algoritmi più sofisticati o sull'ottimizzazione dei parametri, questo lavoro ha indagato se l'arricchimento del dataset attraverso parafrasi possa rendere la dimenticanza più robusta e profonda.

I risultati ottenuti confermano l'ipotesi iniziale: esporre il modello a diverse formulazioni linguistiche della stessa informazione semantica non rappresenta un semplice aumento quantitativo dei dati, ma interviene in modo strutturale sul processo di eliminazione della conoscenza. L'analisi condotta su sei metodi state-of-the-art ha mostrato che l'introduzione di parafrasi produce benefici consistenti sulla riduzione della memorizzazione, in particolare per le tecniche basate su ottimizzazione di preferenze. La diversità sintattica ostacola la rigenerazione letterale delle sequenze memorizzate e favorisce un forgetting più profondo, suggerendo che l'informazione viene rimossa a un livello più astratto rispetto alla semplice cancellazione dei pattern testuali superficiali.

Emerge inoltre una relazione differenziata tra le caratteristiche intrinseche di ciascun metodo e la sua capacità di sfruttare la variabilità linguistica. Questo indica che le parafrasi non costituiscono una soluzione universale, ma uno strumento che amplifica l'efficacia di approcci già sensibili alla struttura semantica dell'informazione. La comprensione di questa dipendenza apre la strada a strategie di unlearning più consapevoli, in cui la scelta del metodo e la progettazione del dataset vengano calibrate congiuntamente.

Nonostante i risultati incoraggianti, lo studio presenta alcune limitazioni che aprono naturalmente a sviluppi futuri. In primo luogo, gli esperimenti sono stati condotti su un singolo modello di dimensioni relativamente contenute (Llama-3.2-1B-Instruct) e su un benchmark sintetico controllato (TOFU). Sebbene questa scelta abbia permesso di isolare con precisione l'effetto delle parafrasi, resta da verificare se gli stessi benefici si mantengano su modelli di scala maggiore e su dataset realistici contenenti conoscenze fattuali complesse o informazioni sensibili reali. Inoltre, la generazione delle parafrasi è stata effettuata attraverso un modello linguistico esterno, introducendo un costo computazionale aggiuntivo nella fase di preparazione dei dati. Approcci alternativi, basati su tecniche di data augmentation più leggere o su procedure di parafrasamento automatico integrate nel processo di training, potrebbero ridurre tale overhead mantenendo i benefici osservati.

Un altro aspetto meritevole di approfondimento riguarda la qualità delle parafrasi stesse. In questo lavoro non è stata condotta un'analisi sistematica sulla diversità sintattica o semantica delle riformulazioni generate, né sul loro grado di equivalenza informativa rispetto alle domande originali. È possibile che parafrasi di qualità inferiore, troppo simili tra loro o eccessivamente distanti dal contenuto originale, possano ridurre l'efficacia del processo o introdurre distorsioni indesiderate. Studi futuri potrebbero investigare criteri di selezione o valutazione della qualità delle parafrasi, al fine di massimizzarne l'impatto sull'unlearning senza compromettere la coerenza semantica del dataset.

Resta inoltre aperta la questione dell'interpretabilità del processo. Sebbene i risultati mostrino che le parafrasi migliorano l'unlearning, i meccanismi attraverso cui ciò avviene a livello di parametri del modello rimangono in gran parte opachi. Comprendere quali strati della rete, quali pattern di attenzione o quali regioni dello spazio di embedding siano maggiormente influenzate dalla variabilità linguistica potrebbe fornire indicazioni preziose per progettare tecniche di unlearning più mirate ed efficienti.

Sul piano delle applicazioni, i risultati di questo lavoro hanno implicazioni concrete per il design di sistemi di unlearning in contesti reali. La constatazione che un numero relativamente contenuto di parafrasi (5-10) sia sufficiente a ottenere miglioramenti significativi rende questa strategia praticabile anche in scenari con risorse computazionali limitate. Aziende e organizzazioni che devono conformarsi a normative sulla privacy, come il GDPR, potrebbero integrare la generazione automatica di parafrasi nei propri workflow di unlearning, riducendo il rischio di persistenza della memorizzazione senza dover ricorrere a riaddestramenti completi del modello.

In conclusione, questo lavoro dimostra che l'introduzione di parafrasi durante il processo di unlearning rappresenta una strategia efficace per rafforzare la rimozione della conoscenza nei modelli di linguaggio. La variabilità linguistica ostacola sia la memorizzazione sintattica che quella semantica, e se dosata in modo appropriato, preserva l'utilità generale del modello.
